\documentclass[../main.tex]{subfiles}
\begin{document}

\FloatBarrier
\section{Summary}
\label{sec:summary}

This thesis studied ordinal recognition of student engagement on a four-level scale, using the features OpenFace extracts from each frame and the features I3D extracts from the whole clip. It set out to improve recognition over standard single-encoder models and to describe the problem accurately, in particular its strong, uneven class imbalance and its behaviour under distribution shift. The method is a complete, reproducible pipeline: feature extraction and leak-free preprocessing; the proposed hybrid model, which splits the 709-dimensional OpenFace features into five groups (gaze, eye landmarks, face landmarks, head pose, action units), encodes each group with its own temporal encoder, and optionally adds a global I3D motion stream while training every stream with a small auxiliary head; five single-encoder baselines; three loss functions; and a 3$\times$3 cross-dataset evaluation protocol. The central idea was tested by training all $3^5=243$ per-group encoder choices, with the I3D stream enabled and disabled, and by a self-collected, test-only private set of 366 hand-labelled clips that serves as the deciding deployment test. Each of these pieces was chosen so that the two goals, better recognition and an honest account of the problem, could be examined with the same evidence rather than argued separately.

The proposed hybrid model with the I3D stream enabled is the best in-domain model and, as a family, outperforms the strongest baseline. The I3D-stream-enabled configurations are centred above the baseline QWK of 0.537, with a median of 0.553 and 82\% of the 243 clearing the baseline, and the best configuration reaches QWK 0.605 with a mixed TCN/Transformer/LSTM assignment. Reporting the family rather than the single best cell is deliberate: because the improvement holds across most of the grid, it reflects the decomposition itself rather than one lucky assignment that might not survive a change of seed or dataset. The gain over the best baseline is nonetheless small, at $+0.068$ QWK, so the value lies in the framework and the full evidence rather than in a jump in state of the art. Two findings explain why the framework is still worth building. First, the fusion was measured before it was built: the clip-level agreement analysis of Appendix~\ref{chap:baseline} showed that the best OpenFace model and the I3D baseline disagree with each other more than OpenFace models disagree among themselves, and that each family classifies clips correctly that the other misses, which the paired I3D ablation then confirmed in-domain with a mean gain of $+0.060$ QWK on the seen-target cells. Second, most groups are not affected by which encoder they get: only head pose shows a clear preference, for a TCN, with a 0.022 QWK spread across the three encoders against $\le0.006$ for the other four groups. The design space is therefore flat, a TCN is a safe default for every group, and the gain comes from splitting and fusing the feature groups rather than from searching for a per-group encoder, which makes the design rule far easier to transfer than a single tuned configuration.

Engagement models do not generalise across datasets, and this is the study's clearest negative result. Off-diagonal QWK drops to near zero, falling to 0.02 from CMOSE to DAiSEE, even when accuracy appears acceptable, so a model that looks competent on its own test set can be worthless on data drawn from a different source. The in-domain leaderboard is of little help in anticipating this, since it predicts transfer only weakly, with a rank correlation of 0.31 for the baselines and 0.27 for the proposed model; a model cannot be selected for deployment from its in-domain score alone. The consequence is a change in how such systems should be measured: ordinal agreement metrics, not accuracy, must be reported on imbalanced engagement data, because accuracy stays high and misleading exactly where it matters least, and training on the combined datasets is the simplest tool that helps when the deployment target is unknown. The self-collected, test-only, in-the-wild private set is where these threads combine into a single deployment result. The proposed model outperforms the best baseline in all nine train$\times$test cells of the cross-dataset matrix; because the private set is test-only, the only remaining choice is the training data, and training on both datasets together with the proposed hybrid model gives the best result at QWK 0.379, outperforming the best baseline by $+0.094$ against $+0.068$ in-domain and surpassing every single-dataset-trained proposed model. That the margin grows from in-domain to the private set is the important point: the same decomposition that is a small in-domain gain becomes a representation that degrades more slowly under distribution shift, which, for a task whose only real test is deployment, is what matters most. Against this stands the one problem the work does not solve. The rare E0 (Highly\nobreakdash-Disengage) class, around 2--3\% of the data, is traded away rather than solved by the proposed model in-domain and drops to zero recall under distribution shift, so the class that a monitoring system most needs to catch is also the one it detects least reliably.

Several limits should be read alongside these conclusions, and each points to work still to be done. The private set was labelled by a single annotator under the public-dataset rubric, so no inter-annotator agreement can be reported, and its 366 clips limit how reliable the private-set conclusions are, especially for the ten E0 clips, where a single wrong clip strongly affects the recall and the collapse to zero should be read as directional rather than exact. On the architecture itself, most of the 243-configuration grid is not affected by which encoder it gets and the best in-domain gain over the baseline is small, so the decomposition should be presented as a sound, interpretable framework rather than a large jump in performance. The weak in-domain signal on DAiSEE, at QWK 0.166, meant that dataset could not be used for architecture development because its ordinal signal is too faint to separate architecture differences from label noise, so the cross-dataset conclusions rest mainly on CMOSE and the private set. Finally, a single seed, learning rate, and encoder size were used throughout, so the reported gaps are point estimates rather than confidence intervals, and small differences between nearby configurations should not be over-interpreted.

\FloatBarrier
\section{Suggestion for Future Work}
\label{sec:future}

The most valuable next step is to target the rarest class directly, since it is the one place where the current model clearly falls short. Combining the architecture with fixes aimed at the imbalance, such as ordinal losses tuned for E0, minority oversampling with SMOTE~\cite{smote}, or cost-sensitive decision thresholds, could recover E0 recall without trading away the E3 performance that the in-domain results depend on. The loss study already hints that this is a trade-off rather than a solved problem: cross-entropy gives the best ordinal agreement while weighted and ordinal losses recover minority recall and balanced ordinal error, with no single loss winning on every count. A method that keeps the ordinal-agreement strength of cross-entropy while borrowing the minority-class sensitivity of the weighted and ordinal losses is therefore the natural target. Because E0 is so rare, even small gains here would change the practical value of the system more than further QWK gains on the common classes, and a deployed system would still be wise to treat low-confidence disengagement predictions with care and, where possible, leave the final call to~a~human.

The cross-dataset gap calls for a dedicated response beyond the plain dataset combination that already helps. Explicit unsupervised domain adaptation or domain-invariant representation learning would attack the transfer collapse head-on rather than diluting it, and a small amount of labelled target data could calibrate the model before deployment where a fully unsupervised approach falls short. Part of the difficulty is that the I3D stream, so useful in-domain, is neutral or slightly harmful under shift because it carries information that is discriminative but specific to each dataset, so a domain-adaptation method that separates the shared engagement signal from dataset-specific motion cues would address both the transfer gap and the fragility of the fusion at once. A related refinement concerns how the group encoders are chosen: the discrete 243-configuration search could be replaced with a differentiable architecture search or a gating mechanism that picks per-group temporal operators end-to-end, so the model learns which encoder suits which group on its own rather than selecting from a listed grid, and could in principle adapt that choice per domain.

Stronger evaluation would tighten every claim made here and is a prerequisite for turning the study's directional findings into settled ones. Multi-annotator labelling of a larger private set, with reported inter-annotator agreement, would put the deployment conclusion on firmer ground and give the rare E0 class enough examples for its recall to be estimated with confidence rather than from ten clips. Multi-seed training with confidence intervals would, in turn, replace the point estimates used throughout with statistically grounded claims, making it possible to say which of the small in-domain differences between nearby configurations are real and which fall within noise.

\end{document}
